\begin{abstract}
Reliable communication over a noisy channel requires redundancy, but adding redundancy lowers the communication rate. For the discrete memoryless channels considered in this paper, Shannon's noisy-channel coding theorem resolves this tension by identifying a capacity: at rates below capacity, families of increasingly long codes can make block-decoding errors arbitrarily unlikely, whereas rates above capacity cannot support reliable communication. The theorem establishes what is possible, but it does not itself provide an explicit and efficient code. This paper introduces polar codes as one structured response to that construction problem. Beginning with a repetition-code example, we develop the communication model, the binary symmetric channel, rate, Hamming distance, and linear codes. We then use typical noise patterns and binary entropy to motivate the capacity of the binary symmetric channel and introduce mutual information for more general channels. Finally, an information-theoretic preview and two- and four-bit binary erasure channel examples show how a recursive transform creates synthetic channels of unequal quality. Placing information in reliable positions and freezing the remaining positions to known values produces a polar code that can be encoded and decoded efficiently. Together, Shannon's theorem and polar codes illustrate the relationship between a theoretical communication limit and one explicit method for approaching it.
\end{abstract}
